\subsection{李代数的包络双代数}

回忆 g 表示李代数，U 表示其包络代数，连同其典范滤过 \((\mathrm{U}_{n})_{n\geqslant0}\)。

命题 7. 在代数 U 上存在且仅存在一个余积 c，使 U 成为双代数且 \(\sigma(g)\) 的元素都是本原的。双代数 (U, c) 是余交换的；其余单位是线性形式 \(\varepsilon\)，使得 U 的每个元素 \(x\) 的常数项（第 I 章，第 2 节，第 1 号）为 \(\varepsilon(x)\) .1。U 的典范滤过 \((\mathrm{U}_n)_{n\geqslant 0}\) 与这一双代数结构相容。

\begin{enumerate}
\def\labelenumi{(\alph{enumi})}
\tightlist
\item
  设 \(x \in \mathfrak{g}\)；我们记 \(c_0(x) = \sigma(x) \otimes 1 + 1 \otimes \sigma(x) \in \mathrm{U} \otimes \mathrm{U}\)。若 \(x, y\) 属于 \(\mathfrak{g}\)，则
\end{enumerate}

\[
c _ {0} (x) c _ {0} (y) = (\sigma (x) \sigma (y)) \otimes 1 + 1 \otimes (\sigma (x) \sigma (y)) + \sigma (x) \otimes \sigma (y) + \sigma (y) \otimes \sigma (x),
\]

由此

\[
[ c _ {0} (x), c _ {0} (y) ] = c _ {0} ([ x, y ]).
\]

由 U 的泛性质（第 I 章，第 2 节，第 1 号，命题 1），存在且仅存在一个幺代数同态

\[
c \colon \mathrm{U} \to \mathrm{U} \otimes \mathrm{U}
\]

使得 \(c(\sigma(x)) = \sigma(x) \otimes 1 + 1 \otimes \sigma(x)\) 对 \(x \in \mathfrak{g}\) 成立。这就证明了命题 7 中的唯一性断言。

\begin{enumerate}
\def\labelenumi{(\alph{enumi})}
\setcounter{enumi}{1}
\tightlist
\item
  我们证明 \(c\) 是余结合的。线性映射 \(c'\) 与 \(c''\)（\(U\) 到 \(U \otimes U \otimes U\) 的、由下式定义的）
\end{enumerate}

\[
c ^ {\prime} = (c \otimes \mathrm{Id} _ {\mathrm{U}}) \circ c \quad \text { and } \quad c ^ {\prime \prime} = (\mathrm{Id} _ {\mathrm{U}} \otimes c) \circ c
\]

是幺代数同态，且在 \(\sigma (\mathfrak{g})\) 上一致，因为对 \(a\in \sigma (\mathfrak{g})\)

\[
c ^ {\prime} (a) = a \otimes 1 \otimes 1 + 1 \otimes a \otimes 1 + 1 \otimes 1 \otimes a = c ^ {\prime \prime} (a),
\]

由此得结论。

\begin{enumerate}
\def\labelenumi{(\alph{enumi})}
\setcounter{enumi}{2}
\item
  我们证明 \(c\) 是余交换的。设 \(\tau\) 为 \(\mathbf{U} \otimes \mathbf{U}\) 的自同构，使得 \(\tau(a \otimes b) = b \otimes a\) 对 \(a, b\) ∈ \(\mathbf{U}\) 成立。映射 \(\tau \circ c\) 与 \(c\)（都是 \(\mathbf{U}\) 到 \(\mathbf{U} \otimes \mathbf{U}\) 的映射）是幺代数同态，且在 \(\sigma(\mathfrak{g})\) 上一致，由此得结论。
\item
  我们证明 \(\varepsilon\) 是 \(c\) 的余单位。U 到 U 的映射 \((\mathrm{Id}_{\mathbb{U}} \otimes \varepsilon) \circ c\) 与 \((\varepsilon \otimes \mathrm{Id}_{\mathbb{U}}) \circ c\) 是幺代数同态，且与 \(\mathrm{Id}_{\mathbb{U}}\) 在 \(\sigma(\mathfrak{g})\) 上一致。
\item
  我们知道 \(\mathrm{U}_0 = \mathrm{K}.1, \mathrm{U}_n \subset \mathrm{U}_{n+1}, \mathrm{U} = \bigcup_{n \geqslant 0} \mathrm{U}_n\) 以及 \(\mathrm{U}_n. \mathrm{U}_m \subset \mathrm{U}_{n+m}\)（第 I 章，第 2 节，第 6 号）。设 \(a_1, \ldots, a_n\) 属于 \(\sigma(\mathfrak{g})\)。那么
\end{enumerate}

\[
\begin{array}{l} c (a _ {1} \ldots a _ {n}) = \prod_ {i = 1} ^ {n} c (a _ {i}) = \prod_ {i = 1} ^ {n} (a _ {i} \otimes 1 + 1 \otimes a _ {i}) \\ = \sum_ {i = 0} ^ {n} \sum_ {\alpha \in I (i)} (a _ {\alpha (1)} \ldots a _ {\alpha (i)}) \otimes (a _ {\alpha (i + 1)} \ldots a _ {\alpha (n)}), \end{array}\tag{8}
\]

其中 \(I(i)\) 表示 \([1, n]\) 的、在每个区间 \([1, i]\) 与 \([i + 1, n]\) 上都递增的置换的集合。由于 \(U_{n}\) 是由 \(\sigma(g)\) 中至多 n 个元素的乘积生成的 K-模，公式 (8) 蕴含滤过 \((\mathrm{U}_{n})\) 与 \((\mathrm{U}, c)\) 的双代数结构相容。

定义 3. 双代数 \((\mathrm{U}, c)\) 称为李代数 \(\mathfrak{g}\) 的包络双代数。

命题 8. 设 E 为双代数，其余积记作 \(c_{\mathbb{E}}\)，并设 h 为 g 到 P(E) 的李代数同态（第 2 号，命题 4）。使对一切 x ∈ g 有 f(σ(x)) = h(x) 的幺代数同态 f: U → E 是双代数态射。

\[
(f \otimes f) \circ c = c _ {\mathrm{E}} \circ f.
\]

\[
a \in \sigma (\mathfrak {g})
\]

\[
(f \otimes f) (c (a)) = f (a) \otimes 1 + 1 \otimes f (a) = c _ {\mathbb {E}} (f (a))
\]

因为 \(f(a) \in \mathrm{P}(\mathrm{E})\)。类似地，若 \(\varepsilon_{\mathbb{E}}\) 是 \(\mathrm{E}\) 的余单位，则 \(\varepsilon_{\mathbb{E}} \circ f\) 是 \(\mathrm{U} \to \mathrm{K}\) 的幺代数同态，它在 \(\sigma(\mathfrak{g})\) 上为零（第 1 号，命题 1），因而与 \(\varepsilon\) 一致。

由命题 5 与命题 8 可知，映射 \(f \mapsto f \circ \sigma\) 给出双代数同态 \(U(\mathfrak{g}) \to E\) 与李代数同态 \(\mathfrak{g} \to P(E)\) 之间的一一对应。

推论. 设 \(\mathfrak{g}_i\)（\(i = 1,2\)）为李代数，\(\mathrm{U}(\mathfrak{g}_i)\) 为其包络双代数，\(\sigma_i\colon \mathfrak{g}_i\to \mathrm{U}(\mathfrak{g}_i)\) 为典范映射。对每个李代数同态 \(h\colon \mathfrak{g}_1\to \mathfrak{g}_2\)，幺代数同态 \(\mathrm{U}(h)\colon \mathrm{U}(\mathfrak{g}_1)\to \mathrm{U}(\mathfrak{g}_2)\)（它使 \(\mathrm{U}(h)\circ \sigma_1 = \sigma_2\circ h\)，第 I 章，第 2 节，第 1 号）是双代数态射。

